\chapter{Исследование и построение решения}\label{ch:ch2}

\section{Построение частично упорядоченного множества операций с памятью в многопоточной программе с учётом модели памяти}\label{sec:ch2/sec1}

В граве RISC-V Weak Memory Order~\ref{sec:ch1/rvvmo} бал определён програмный порядок во время исполнения в одном потоке.
Но исследование такого порядка в изоляции не имеет большого интереса, так как он формирует полный порядок~\fixme{total order}.

Неопределённость относительного порядка исполнения комманд является отличительной особенностью многопоточной программы.
Интерес для иследования представляет глобальный порядок операций в памяти (GMO).
В частности -- определение, что все строгие программные порядки (PPO) сохранены в GMO.

Пусть $Prog$ -- многопоточная программа, заданная следующим образом:

\begin{enumerate}
        \item программа состоит из $n$ функций $func$, исполняемых одновременно в разных потоках;
        \item функция состоит из мгновенных~\fixme{здесь хочется сказать атомарных} операций с памятью и арифметических инструкций;
        \item мгновенные операции с памятью могут пересекаться по адресам с операциями в других функциях;
\end{enumerate}

Определим \textbf{операцию с памятью} как семантически неделимую структуру, состоящую из мгновенных операций с памятью или из барьера.
Например, набор инструкций из листинга~\ref{lst:cas} содержит в себе две мгновенные операции с памятью Lr и Sc, но семантически отражают одну операцию AMOCAS~\ref{sec:ch1/zacas}.
Если операцию AMOCAS в цикл, как на листинге~\ref{lst:amocas-loop}, то семантически получится операция AMOSWAP~\fixme{check asm}.

\begin{minipage}{\textwidth}
\centering
\lstinputlisting[caption={AMOCAS в цикле.}, label={lst:amocas-loop}, captionpos=b]{src/amocas-loop.s}
\end{minipage}

Таким образом, определим 3~\fixme{do not need fence?} семантические категории операций с памятью:

\begin{enumerate}
        \item Чтение(R) -- атомарное чтение из памяти;
        \item Запись(W) -- атомарная запись в память;
        \item Чтение-Модификация-Запись(RMW) -- атомарная модификация память;
\end{enumerate}

Отдельные Lr,Sc или AMOCAS не могут быть выраженны в этих категориях, но если их повторять в цикле, пока результат операций не будет успешен, то такие структуры уже являются RMW операциями с памятью.

Построим частично упорядоченное множество $P$ из многопоточной программы по формуле~\ref{eq:prog2poset}.

\begin{equation}
    \label{eq:prog2poset}
    \begin{multlined}
      P : \forall x \in P \equiv x \text{ -- операция с памятью} \\
      <_p : \forall x, y \in P \mid x <_p y \equiv xy \text{ -- PPO}
    \end{multlined}
\end{equation}

Линейный порядок элементов множества $P$ определяет GMO -- некоторый вариант исполнения программы.
Верификация конкретного результата исполнения программы заключается в поиске GMO, который даст такой же результат~\fixme{reformulate}.
Если такого GMO не существует, то в тестируемой системе обнаружена ошибка.
Главной задачей будет построение всех линейных порядков частично упорядоченного множества, которое определяется многопоточной программой и моделью памяти системы, где эта программа исполняется.


\section{Топологическая сортировка}\label{sec:ch2/topol}

Обозначим $l_P$ -- линейный порядок множества $P$, $L(P) : \forall l_P \in L(P)$.
Линейный порядок множества $P$ тоже самое, что типологическая сортировка графа $G$, поэтому $L(P) \equiv L(G) \equiv L$

Для построения $L(P)$ используется алгоритм топологической сортировки~\ref{alg:topo} аналогично формуле~\ref{eq:lin-ext}.

\begin{algorithm}[H]
\caption{Обход частично упорядоченного множества топологической сортировкой}
\label{alg:topo}
\begin{algorithmic}[1]
\Require $P$ -- частично упорядоченное множество. $l_P$ -- текущий линейный порядок.
\Function{Traverse}{$P$, $l_P$}
\If{$P = \emptyset$}
  \State $L(P) = l_P \cup L(P)$ \Comment Полностью построили линейный порядок
\EndIf
\ForAll{$\forall x \in min(P)$}
  \State \Call{Process}{$x$} \Comment Обработка элемента множества
  \State \Call{Traverse}{$P \setminus x$, $l_P \stackrel{add}{\longleftarrow} x$}
\EndFor
\EndFunction
\end{algorithmic}
\end{algorithm}

Расчитаем количесвтво линейных порядков в $L(P)$ и сложность алгоритма их построения. Частично упорядоченое множество может быть описано ацикличным направленным графом.
Рассмотрим случайный DAG -- $G$, построенный следующим образом:
\begin{itemize}
  \item фиксировано число вершин $n = |V(G)|$;
  \item $\sigma_i \in V(G) \mid 1 \leq i \leq n$;
  \item выбираем любую перестановку из $\sigma_1 \dots \sigma_n$. Это будет заданный топологический порядок DAG;
  \item Независимо $\forall (i, j) \mid i < j$, проводим ребро $(\sigma_i, \sigma_j) \in E(G)$ с вероятностью $p$;
\end{itemize}
Данное построение, гарантирует построение корректного DAG и позволяет варировать его топологию параметром $p$.

\subsection{Количство линейных порядков}\label{sec:ch2/linext-count}

Обозначим $\sigma_i$ -- вершина $\sigma \in V(G)$, с порядковым номером $i$. И $\text{id}^{\alpha}(\sigma_i)$ -- позиция вершины $\sigma_i$ в топологическом порядке $\alpha$.
Пусть $\tau$ -- топологический порядок, по которому строился граф и $\text{id}^{\tau}(\sigma_i) \equiv i$.
Тогда $\forall (\sigma_i, \sigma_j) \in E(G) | \mid i < j$ по построению.

Возьмём другой порядок $\pi$. $\pi$ является топологическим порядком графа $G$, тогда и только тогда, если в графе нет ребра, направленного в противоположную сторону относительно $\pi$.
\begin{equation}
  \nexists (\sigma_i, \sigma_j) \in E(G) \mid \text{id}^{\pi}(\sigma_j) < \text{id}^{\pi}(\sigma_i)
\end{equation}

Для фиксированного порядка $\pi$, количество пар вершин $(\sigma_i, \sigma_j) | i < j$, но расположенных в $\pi$ в обратном порядке $\text{id}^{\pi}(\sigma_j) < \text{id}^{\pi}(\sigma_i)$ в точности равно $\text{inv}(\pi)$ -- числу инверсий перестановки $\pi$ относительно $\tau$.\fixme{find cite}
Чтобы перестановка $\pi$ была корректной топологической сортировкой $G$, ниодного потенциально неправильно направленного ребра не должно существовать.

Вероятность, что ребра не существует -- $1-p$ и рёбра расставляются независимо, тогда:
\begin{equation}
  \mathbb{P}[\pi \text{ -- топологическая сортировка}] = (1-p)^{\text{inv}(\pi)}
\end{equation}
а общее количество линейных порядков суммируется по всем возможным $n!$ перестановкам \fixme{find cite}:
\begin{equation}
  \mathbb{E}(|L|) = \sum_{\pi}{(1-p)^{\text{inv}(\pi)}} = \prod_{k=1}^{n}{\frac{1-q^k}{1-q}}, \text{ где } q = 1-p
\end{equation}

Оценим некоторые пределы.
\begin{subequations}
    \label{eq:linext-count-lims}
    \begin{align}
      p & = 1,\ \mathbb{E} = 1 \label{eq:linext-count-lims-p1} \\
      p & = 0,\ \mathbb{E} = \sum_{\pi}{1^{\text{inv}(\pi)}} = n! \label{eq:linext-count-lims-p0} \\
      p & = \frac{1}{2},\ \mathbb{E} = 2^n \prod_{k=1}^{n} (1 - \frac{1}{2}^k) \rightarrow 0.29*2^n, \text{ при } n \rightarrow \infty \label{eq:linext-count-lims-p12}
    \end{align}
\end{subequations}
$Q_n = \prod_{k=1}^{n} (1 - \frac{1}{2}^k), Q_{\infty} \approx 0.29$.
Случай \cref{eq:linext-count-lims-p1} соответсвует единственному допустимому линейному порядку, по которому строился DAG.

\subsection{Количство минимальных элементов}\label{sec:ch2/min-elems}

Расчитаем, среднее количество минимальных элементов в случайном DAG.

Выберем и зафиксируем случайную вершину $\sigma$ и выберем топологическую перестановку, использованную при построении графа $\tau$.
Пусть $\text{id}^{\tau}(\sigma) = k$ -- позиция $\sigma$ с вероятностью:
\begin{equation}
  \label{eq:p-pos-k}
  \mathbb{P}[\text{id}^{\tau}(\sigma) = k] = \frac{1}{n}
\end{equation}
Тогда ребро в $\sigma$ может входить только в из вершин, стоящих левее в порядке с вероятностью:
\begin{equation}
  \label{eq:no-in-k}
  \mathbb{P}[\text{нет входящих рёбер} \mid \text{id}^{\tau}(\sigma) = k] = (1-p)^{k-1}
\end{equation}
И из формул~\ref{eq:p-pos-k}~\ref{eq:no-in-k} следует вероятность $\sigma \in \text{min}(G)$:
\begin{equation}
  \label{eq:n-is-min}
    \begin{multlined}
      \mathbb{P}[\sigma \in \text{min}(G)] = \\
      = \sum_{k=1}^{n}{\mathbb{P}[\text{id}^{\tau}(\sigma) = k]\mathbb{P}[\text{нет входящих рёбер} \mid \text{id}^{\tau}(\sigma) = k]} = \\
      = \sum_{k=1}^{n}{\frac{1}{n}(1-p)^{k-1}} = \frac{1}{n}\frac{1-q^n}{1-q}, \text{ где } q = 1-p
    \end{multlined}
\end{equation}
И математическое ожидание количества минимальных элеметов:
\begin{equation}
  \label{eq:num-of-min}
  \mathbb{E}(|min(G)|) = \frac{1-q^n}{1-q}
\end{equation}

\subsection{Сложность алгоритма построения всех линейных порядков алгоритмом~\ref{alg:topo}}\label{sec:ch2/topoalg-comp}

Рассмотрим $N$ шаг рекурсии в алгоритме ~\ref{alg:topo}, где $N = |P|$.
Вычислим сколько операций $C_N$ занимает вычисление одного шага рекурсии.
Вычисление минимальных элементов требует $|P|$ операций, а одна итерация цикла занимает $C_{N-1}$ операций.
Вместе с формулой~\cref{eq:num-of-min} Получаем рукурсивную формулу вычисления количества операций на $N$-ом шаге рекурсии:
\begin{equation}
  \begin{multlined}
    \label{eq:topo-comp}
    C_N = |P| + \mathbb{E}(|min(G)|)C_{N-1} = N + \frac{1-q^N}{1-q}C_{N-1} = \dots = \\
    = \sum_{k=1}^{N}{k\prod_{l=1}^{N-k}{\frac{1-q^{k+l}}{1-q}}}
  \end{multlined}
\end{equation}

Оценим некоторые пределы.
\begin{subequations}
    \label{eq:topo-comp-lims}
    \begin{align}
      q & \rightarrow 0,\ C_N \approx \frac{N(N-1)}{2} \label{topo-comp-lims-p1} \\
      q & = 1 - \varepsilon,\ \varepsilon \rightarrow 0, \ C_N \approx N! \sum_{k=1}^{N}{\frac{1}{(k-1)!}} \rightarrow eN!, \text{ при } N \rightarrow \infty \label{topo-comp-lims-p1} \\
      q & = \frac{1}{2},\ C_N = \sum_{k=1}^{N}{k 2^{N-k}\prod_{j=k+1}^{N}{1-2^{-j}}} = \\
        & = \sum_{k=1}^{N}{k 2^{N-k} \frac{Q_N}{Q_k}} \approx \\
        & \approx Q_{\infty}\sum_{m=0}^{N-1}\frac{2^m}{Q_{N-m}} \approx \\
        & \lgroup Q_{N-m} \approx Q_{\infty}, \text{ при } m \ll N \rgroup \\
        & \approx \sum_{m=0}^{N-1}{(N-m)2^m} = 2^{N+1} - N - 2 \label{eq:topo-comp-lims-p12}
    \end{align}
\end{subequations}
